% =====================================================================
%  Chapitre 2 — Analyse des besoins et conception
% =====================================================================
\chapter{Analyse des besoins et conception}
\label{chap:analyse}

\section*{Introduction}
Après avoir situé le projet, ce chapitre en établit les fondations. Nous
identifions d'abord les acteurs de la plateforme et exprimons les besoins
fonctionnels et non fonctionnels. Nous présentons ensuite la conception à
travers le diagramme de cas d'utilisation global, le diagramme de classes
et l'architecture applicative. Nous terminons par l'initialisation du
projet : le backlog produit, la planification des sprints et
l'environnement de développement.

\section{Analyse des besoins}

\subsection{Identification des acteurs}
Un acteur désigne un rôle joué par une entité externe qui interagit avec le
système. Pour Winaity Template Builder, nous distinguons les acteurs
suivants.

\begin{itemize}[leftmargin=1.4cm]
  \item \textbf{Utilisateur / Éditeur.} Membre d'une entreprise cliente. Il
  crée, modifie, prévisualise et exporte des modèles, utilise l'assistant
  d'IA et organise ses modèles (favoris, duplication).

  \item \textbf{Administrateur d'entreprise (Admin).} Utilisateur disposant
  de droits étendus au sein de son entreprise. Outre les capacités d'un
  éditeur, il invite et gère les membres de son équipe, souscrit un
  abonnement et gère la facturation, et génère des clés d'intégration
  (API).

  \item \textbf{Membre marketing.} Rôle spécialisé se comportant comme un
  éditeur, mais habilité en plus à constituer et maintenir la galerie de
  modèles prédéfinis propre à son entreprise, par catégorie.

  \item \textbf{Super-administrateur.} Acteur interne à l'éditeur de la
  plateforme. Il supervise l'ensemble des entreprises, consulte les
  abonnements et attribue les plans, y compris le plan interne réservé aux
  entités du groupe.

  \item \textbf{Client applicatif (M2M).} Système externe (par exemple un
  CRM) authentifié par une clé d'API, qui embarque l'éditeur pour ses
  propres utilisateurs via le flux d'intégration.
\end{itemize}

\subsection{Besoins fonctionnels}
Les besoins fonctionnels décrivent les services que la plateforme doit
rendre. Ils sont regroupés ci-dessous par domaine.

\paragraph{Authentification et gestion des comptes.}
\begin{itemize}[leftmargin=1.4cm,nosep]
  \item S'inscrire et créer une entreprise, ou rejoindre une entreprise
  existante sur invitation.
  \item S'authentifier de manière sécurisée (jeton JWT) et gérer sa
  session.
  \item Consulter et modifier son profil.
\end{itemize}

\paragraph{Multi-entreprises et gestion des utilisateurs.}
\begin{itemize}[leftmargin=1.4cm,nosep]
  \item Garantir l'isolation stricte des données par entreprise.
  \item Inviter des membres, accepter une invitation, lister l'équipe.
  \item Attribuer des rôles (administrateur, éditeur, marketing).
\end{itemize}

\paragraph{Édition de modèles multi-canal.}
\begin{itemize}[leftmargin=1.4cm,nosep]
  \item Créer, modifier, dupliquer, supprimer et lister des modèles.
  \item Éditer un e-mail par glisser-déposer de blocs, prévisualiser en
  direct et éditer le code.
  \item Créer des factures, contrats, SMS et messages RCS.
  \item Exporter en HTML responsive et en PDF, substituer des variables de
  personnalisation.
  \item Gérer les favoris et la galerie de modèles prédéfinis.
\end{itemize}

\paragraph{Assistance par intelligence artificielle.}
\begin{itemize}[leftmargin=1.4cm,nosep]
  \item Générer un e-mail complet à partir d'une consigne en langage
  naturel.
  \item Modifier un modèle en conversation (thème, image, disposition,
  ajout ou déplacement de blocs).
  \item Appliquer une modification à un élément précis sélectionné dans le
  canevas.
  \item Reconstruire un e-mail éditable à partir d'une affiche marketing.
\end{itemize}

\paragraph{Abonnements, intégrations et administration.}
\begin{itemize}[leftmargin=1.4cm,nosep]
  \item Souscrire, changer, annuler et réactiver un abonnement~; appliquer
  les quotas par plan.
  \item Générer des clés d'API et embarquer l'éditeur depuis un outil
  externe.
  \item Superviser les entreprises et attribuer les plans (super-admin).
\end{itemize}

\subsection{Besoins non fonctionnels}
Les besoins non fonctionnels définissent les contraintes de qualité que la
plateforme doit respecter.

\begin{itemize}[leftmargin=1.4cm]
  \item \textbf{Sécurité et confidentialité.} Authentification par jeton,
  cloisonnement rigoureux des données par \texttt{tenant\_id}, contrôle
  d'accès basé sur les rôles, secrets tenus hors du dépôt de code, et
  \textbf{souveraineté des données} garantie par un modèle de langage
  auto-hébergé.
  \item \textbf{Performance.} Génération d'un e-mail par l'IA en une dizaine
  de secondes, rendu et prévisualisation réactifs.
  \item \textbf{Évolutivité (scalabilité).} Architecture microservices
  permettant de faire évoluer et de déployer chaque service
  indépendamment.
  \item \textbf{Maintenabilité.} Séparation des responsabilités, patron
  CQRS côté services, modèle de données partagé entre édition manuelle et
  édition IA.
  \item \textbf{Utilisabilité.} Interface intuitive, accessible à un
  utilisateur métier sans compétence technique.
  \item \textbf{Fiabilité.} Confirmations honnêtes de l'assistant (aucun
  succès annoncé si le modèle n'a pas réellement changé), garde de
  contraste garantissant la lisibilité, gestion propre des erreurs.
\end{itemize}

\section{Conception}

\subsection{Diagramme de cas d'utilisation global}
Le diagramme de cas d'utilisation offre une vue d'ensemble du comportement
fonctionnel de la plateforme et des interactions entre acteurs et cas
d'utilisation. La figure~\ref{fig:uc-global} présente le diagramme global du
système.

% [FIGURE : diagramme de cas d'utilisation global — source cas-utilisation.wsd]
\figplaceholder{Diagramme de cas d'utilisation global}{fig:uc-global}{8cm}

\subsection{Diagramme de classes}
Le diagramme de classes décrit la structure statique du système : les
entités, leurs attributs et leurs relations. Les principales entités du
domaine sont \texttt{Tenant} (entreprise), \texttt{User} (utilisateur et
son rôle), \texttt{Template} (modèle, caractérisé par son type de canal et
son contenu), les entités d'invitation et de jeton, ainsi que les entités
liées à l'abonnement (plan, cycle de facturation, statut). La
figure~\ref{fig:classes} présente le diagramme de classes.

% [FIGURE : diagramme de classes — source diagramme-classes.wsd]
\figplaceholder{Diagramme de classes}{fig:classes}{9cm}

\subsection{Architecture de l'application}
La plateforme adopte une \textbf{architecture microservices}. Cette section
en présente les vues physique et logique.

\subsubsection{Architecture physique}
L'architecture physique décrit les composants déployés et leurs
interactions. Elle se compose des éléments suivants~:

\begin{itemize}[leftmargin=1.4cm]
  \item \textbf{Client (navigateur).} Interface utilisateur exécutée dans le
  navigateur.
  \item \textbf{Application frontend.} Application Next.js servant
  l'interface et hébergeant les routes d'IA côté serveur.
  \item \textbf{Passerelle API (API Gateway).} Unique surface REST publique~;
  elle traduit les appels REST en appels gRPC vers les services internes et
  gère l'authentification, la facturation et l'application des quotas.
  \item \textbf{Service d'authentification.} Gère les comptes, les
  entreprises, les rôles et les abonnements.
  \item \textbf{Service de modèles.} Gère le cycle de vie des modèles
  (CRUD, rendu, duplication, favoris).
  \item \textbf{Services d'IA et d'images.} Services auxiliaires
  (suggestions de palettes, mapping de variables, recherche d'images) et
  \textbf{serveur de modèle de langage auto-hébergé}.
  \item \textbf{Bases de données PostgreSQL} (\texttt{auth\_db} et
  \texttt{templates\_db}) et \textbf{stockage objet MinIO} pour les médias.
\end{itemize}

% [FIGURE : architecture physique — source architecture-technique.drawio]
\figplaceholder{Architecture physique de l'application}{fig:archi-phys}{7cm}

\subsubsection{Architecture logique}
Sur le plan logique, la plateforme suit un modèle en couches communiquant
selon un schéma requête/réponse. Le client émet une requête REST
(authentifiée par un jeton JWT porté par un \emph{cookie}) vers la
passerelle API. Celle-ci relaie la requête, via \textbf{gRPC}, au service
concerné. Les services, développés avec le framework NestJS, appliquent le
patron \textbf{CQRS} (séparation des commandes et des requêtes) et
persistent leurs données au moyen de l'ORM TypeORM. La
figure~\ref{fig:archi-log} illustre cette architecture.

% [FIGURE : architecture logique — source architecture-logique.drawio]
\figplaceholder{Architecture logique de l'application}{fig:archi-log}{7cm}

Le choix du \textbf{gRPC} pour la communication inter-services se justifie
par ses performances (sérialisation binaire via Protocol Buffers) et par le
contrat d'interface fort qu'il impose. Le jeton JWT transporte la charge
utile \texttt{\{ userId, tenantId, email, role \}}, où le
\texttt{tenant\_id} constitue la clé d'isolation présente dans tous les
services.

\section{Initialisation du projet}

\subsection{Backlog produit}
Le backlog produit recense les fonctionnalités attendues, exprimées sous
forme d'\emph{epics} et priorisées. Le tableau~\ref{tab:backlog} en présente
la synthèse.

\begin{table}[H]
\centering
\caption{Backlog produit (synthèse par epic)}
\label{tab:backlog}
\small
\begin{tabularx}{\linewidth}{clXc}
\toprule
\textbf{ID} & \textbf{Epic} & \textbf{Description} & \textbf{Priorité} \\
\midrule
1  & Authentification & Inscription, connexion sécurisée (JWT), profil. & Haute \\
2  & Multi-entreprises & Isolation des données par \texttt{tenant\_id}. & Haute \\
3  & Gestion des utilisateurs & Invitations, équipe, rôles (admin / éditeur / marketing), super-admin. & Haute \\
4  & Éditeur d'e-mail & Canevas de blocs, glisser-déposer, aperçu, code, export HTML/PDF. & Haute \\
5  & Modèles multi-canal & Facture, contrat, SMS, RCS. & Haute \\
6  & Bibliothèque de modèles & Favoris, duplication, galerie de modèles prédéfinis. & Moyenne \\
7  & Génération IA par consigne & E-mail complet à partir d'une invite en langage naturel. & Haute \\
8  & Édition IA conversationnelle & Modification par dialogue et édition ciblée par sélection. & Haute \\
9  & Affiche $\rightarrow$ e-mail & Reconstruction d'un e-mail éditable à partir d'une image. & Haute \\
10 & MCP \& LLM auto-hébergé & Exposition des outils via MCP, inférence sur serveur vLLM privé. & Moyenne \\
11 & Abonnements et facturation & Plans, paiement Stripe, TVA, quotas par plan. & Haute \\
12 & Super-administration & Supervision des entreprises, attribution des plans. & Moyenne \\
13 & Intégrations et embed & Clés d'API, embarquement de l'éditeur par un outil externe. & Moyenne \\
14 & Site marketing & Site public de présentation (accueil, offres, contact\dots). & Basse \\
15 & Mise en production & Conteneurisation, déploiement, exposition sécurisée. & Haute \\
\bottomrule
\end{tabularx}
\end{table}

\subsection{Planification des sprints}
Le projet a été organisé en quatre sprints regroupant les epics du backlog
de manière cohérente, comme le résume le tableau~\ref{tab:sprints}.

\begin{table}[H]
\centering
\caption{Planification des sprints}
\label{tab:sprints}
\small
\begin{tabularx}{\linewidth}{lX}
\toprule
\textbf{Sprint} & \textbf{Contenu} \\
\midrule
Sprint 1 & Authentification, multi-entreprises et gestion des utilisateurs (epics 1--3). \\
Sprint 2 & Éditeur de modèles multi-canal et bibliothèque de modèles (epics 4--6). \\
Sprint 3 & Assistant d'intelligence artificielle agentique (epics 7--10). \\
Sprint 4 & Abonnements, intégrations et mise en production (epics 11--15). \\
\bottomrule
\end{tabularx}
\end{table}

\subsection{Environnement de développement}

\subsubsection{Environnement matériel}
Le développement s'est appuyé sur les ressources suivantes~:
\begin{itemize}[leftmargin=1.4cm,nosep]
  \item un poste de travail personnel~: [À COMPLÉTER — processeur, mémoire
  RAM, système d'exploitation]~;
  \item un \textbf{serveur d'entreprise doté d'un accélérateur graphique
  (GPU)} hébergeant le moteur d'inférence vLLM et servant également au
  déploiement en production.
\end{itemize}

\subsubsection{Environnement logiciel et technologies}
Le tableau~\ref{tab:techno} récapitule la pile technologique. Les
principaux choix sont ensuite justifiés.

\begin{table}[H]
\centering
\caption{Pile technologique du projet}
\label{tab:techno}
\small
\begin{tabularx}{\linewidth}{lX}
\toprule
\textbf{Couche} & \textbf{Technologies} \\
\midrule
Frontend & Next.js~16 (App Router), React~19, TypeScript, Tailwind CSS, Framer Motion \\
Édition & dnd-kit (glisser-déposer), Monaco (éditeur de code), MJML (rendu e-mail), Tiptap (texte riche) \\
Backend & NestJS, gRPC / Protocol Buffers, patron CQRS, TypeORM \\
Données & PostgreSQL (\texttt{auth\_db}, \texttt{templates\_db}), MinIO (stockage objet) \\
IA & Vercel AI SDK, MCP SDK, serveur vLLM auto-hébergé (Gemma / Qwen), services Python auxiliaires \\
Paiement & Stripe (Checkout intégré, abonnements, TVA) \\
Industrialisation & Docker \& Docker Compose, GitLab (CI), Cloudflare Tunnel \\
Outils & Visual Studio Code, Git, Postman, PlantUML / draw.io \\
\bottomrule
\end{tabularx}
\end{table}

\paragraph{Frontend — Next.js et React.}
Next.js, fondé sur React, a été retenu pour son \emph{App Router}, sa
capacité à exécuter du code côté serveur (indispensable pour héberger en
toute sécurité les appels au modèle de langage) et son rendu performant.
React apporte son modèle de composants et son écosystème mature~\cite{nextjs,react}.

% [FIGURE : logos Next.js et React]
\figplaceholder{Next.js et React}{fig:logo-next-react}{2.5cm}

\paragraph{Backend — NestJS et gRPC.}
NestJS structure le code serveur autour de modules, contrôleurs et
services, et facilite l'application du patron CQRS. La communication
inter-services repose sur gRPC, qui offre un contrat d'interface fort et une
sérialisation binaire efficace via Protocol Buffers~\cite{nestjs,grpc}.

% [FIGURE : logos NestJS et gRPC]
\figplaceholder{NestJS et gRPC}{fig:logo-nest-grpc}{2.5cm}

\paragraph{Données — PostgreSQL, TypeORM et MinIO.}
PostgreSQL assure la persistance relationnelle des deux bases du système.
TypeORM fait le pont entre les entités du domaine et la base. MinIO, compatible
avec l'API S3, stocke les médias (images) de la plateforme~\cite{postgresql,typeorm,minio}.

% [FIGURE : logos PostgreSQL et MinIO]
\figplaceholder{PostgreSQL et MinIO}{fig:logo-pg-minio}{2.5cm}

\paragraph{Rendu e-mail — MJML.}
MJML est un langage de balisage qui compile vers du HTML d'e-mail
responsive compatible avec la majorité des clients de messagerie~; il
constitue le format d'export de l'éditeur d'e-mail~\cite{mjml}.

\paragraph{Intelligence artificielle — Vercel AI SDK, MCP et vLLM.}
Le Vercel AI SDK orchestre la boucle agentique d'appel d'outils côté
serveur. Le Model Context Protocol (MCP) expose les outils de construction
au modèle de manière standardisée. L'inférence est assurée par un serveur
\textbf{vLLM auto-hébergé} exécutant le modèle Gemma (et Qwen en repli),
tous deux compatibles avec l'API OpenAI, ce qui garantit que les données ne
quittent pas l'infrastructure du groupe~\cite{vercelaisdk,mcp,vllm,gemma}.

% [FIGURE : logos Vercel AI SDK, vLLM et Gemma]
\figplaceholder{Vercel AI SDK, vLLM et Gemma}{fig:logo-ai}{2.5cm}

\paragraph{Paiement — Stripe.}
Stripe gère les abonnements récurrents, le Checkout intégré et le calcul de
la TVA~\cite{stripe}.

\paragraph{Industrialisation — Docker, GitLab et Cloudflare.}
Docker et Docker Compose conteneurisent l'ensemble des services et
garantissent la reproductibilité des environnements. Le dépôt est hébergé
sur GitLab. L'exposition publique en production s'appuie sur un tunnel
Cloudflare, sécurisé et sans ouverture de port entrant~\cite{docker,cloudflaretunnel}.

\section*{Conclusion}
Ce chapitre a précisé les acteurs et les besoins de la plateforme, présenté
sa conception (cas d'utilisation, classes, architecture microservices) et
posé les bases de sa réalisation : backlog produit, planification en quatre
sprints et environnement technique. Les chapitres suivants détaillent, sprint
par sprint, la mise en \oe{}uvre de ces choix.
